\documentclass{situationist-preprint}

\title{Grounding and Causation: An Analogy Examined}
\msid{MET-14-b}
\msfamily{MET-14}
\msdiscipline{Metaphysics}
\mstype{Companion manuscript}
\mskeywords{grounding; causation; transitivity; metaphysical causation; structural equations; unity of grounding}

\begin{document}
\maketitle

\begin{abstract}
Grounding is often modelled on causation: both are asymmetric relations of dependence that back explanations, and some have proposed that grounding is a kind of causation. MET-14-a argues that grounding is a single relation whose apparent varieties are its restrictions to pairs of levels. We examine whether the analogy with causation supports that thesis. It does in one respect: causation, too, is one relation across many domains, and no one counts chemical and biological causation as distinct relations. But the analogy fails where the unity theorem needs it. Generic grounding is the ancestral of its links and is transitive; causation, on the account of MET-04, is not the ancestral of its links and is not transitive. Grounding is therefore not a species of causation, and the unity of grounding does not rest on the analogy.
\end{abstract}

\section{The analogy}

Three proposals connect grounding and causation. \citet{Schaffer2016} models grounding with the structural equations used for causation, so that grounding becomes a relation of dependence among variables with its own counterfactual structure. \citet{Wilson2018} goes further and proposes that grounding is metaphysical causation: a relation of the same type as causation, distinguished only by the laws that back it. \citet{Bennett2017} places grounding among a family of building relations, alongside composition, constitution and realization, and asks what the family has in common. Each proposal invites the question we address: does the unity of causation support the unity of grounding, and is grounding itself a kind of causation?

The analogy has a clear positive side. Causation operates in physics, chemistry, physiology and psychology, by collisions, reactions, hormonal signals and decisions. No one infers from this diversity of mechanisms that there are as many causal relations as there are sciences. The mechanisms are what carry causal influence at a level; the relation they carry is one. MET-14-a says the same of grounding: logical, set-theoretic, mereological, determinable, realization and normative grounding are one relation applied at different levels (Theorem~3.6 of MET-14-a), and the varieties differ in the pairs of levels at which they operate (their signatures, Definition~2.4 of MET-14-a), as causal mechanisms differ in the domains in which they operate.

\section{Where the analogy fails}

The unity theorem has two premises. Mixed acyclicity (Assumption~2.9 of MET-14-a) says that no chain of links returns to its origin. Level closure (Assumption~2.10) says that whenever a chain of links leads from $f$ to $g$ and the levels of $f$ and $g$ fall in the signature of a variety, $f$ grounds $g$ in that variety. Given both, generic grounding, the relation that holds when some finite chain of links connects two facts (Definition~2.3), is the unique unifier and a strict partial order. In particular it is transitive, and every variety is transitive within its signature.

The definition of generic grounding has a famous causal counterpart. \citet{Lewis1973} defined causation as the ancestral of counterfactual dependence: $c$ causes $e$ if a chain of dependences leads from $c$ to $e$. He did so to make causation transitive, as generic grounding is. The construction failed for causation. MET-04 argues that a cause is a counterfactual difference-maker selected by norms of normality, and it accepts the cases that show causation not to be transitive. A dog bites off a terrorist's right forefinger; he presses the detonator with his left; the bite causes the use of the left hand, which causes the explosion, but the bite does not cause the explosion \citep{McDermott1995,Hall2000}. Each link is causal and the chain is not.

\begin{proposition}\label{p:closure}
The causal counterpart of level closure is false: a chain of causal links between events at given levels does not in general make the first a cause of the last.
\end{proposition}

\begin{proof}
The dog-bite case is a chain of two causal links between events of the same level, physiological and mechanical, whose endpoints are not related as cause and effect. On the account of MET-04-a, the explosion does not depend on the bite when the variables of an apt model are held fixed at their actual values, and the account delivers this verdict.
\end{proof}

The failure has a source that has no counterpart in grounding. Causation, on MET-04, is dependence plus default: whether a factor counts as a cause depends on which variables are held fixed and on which alternatives are normal. A chain of causes can pass through a factor whose effect is cancelled by a default alternative path, as the bite's effect is cancelled by the left hand. Grounding involves no selection by normality and no default paths. If $f$ grounds $g$ and $g$ grounds $h$, there is no alternative route by which $h$ would have obtained without $g$ that cancels $f$'s contribution, since grounding is not a relation of difference-making against a background of defaults but a relation of obtaining in virtue of. This is why Principle~P1 of MET-14-a, that a mixed chain transmits the in-virtue-of relation, is plausible for grounding and its causal analogue is not.

Two further disanalogies point in the same direction. Full grounds necessitate what they ground, with the necessity appropriate to their level, while causes do not necessitate their effects; and grounding is synchronic, while causation, in the cases that matter, runs forward in time. Neither disanalogy is decisive alone. Together with Proposition~\ref{p:closure}, they show that grounding and causation differ in their structural profile and not only in their relata.

\section{Consequences}

\begin{corollary}\label{c:species}
Grounding is not a species of causation.
\end{corollary}

\begin{proof}
If grounding were a kind of causation, as metaphysical causation, it would be a restriction of causation to certain relata. A restriction of an intransitive relation may be transitive, so the inference needs a further step. By MET-14-a, grounding is a single relation including every variety, and its mediate form is transitive across all levels. A relation of the causal type, defined as dependence plus default, is not transitive across its links, by Proposition~\ref{p:closure}; and the default structure that makes it so would, if present in grounding, make mixed chains fail as causal chains fail. Since mixed chains in grounding do not fail (Principle~P1 and Assumption~2.10 of MET-14-a), grounding lacks the feature that makes causation what it is on MET-04.
\end{proof}

Corollary~\ref{c:species} tells against Wilson's identification but not against the structural-equation framework, which is a modelling tool and can represent a transitive relation as easily as an intransitive one. The apparent failures of transitivity in grounding that \citet{Schaffer2012} presents confirm the diagnosis rather than the analogy. As Section~8 of MET-14-a argues, they concern difference-making, which is contrastive: the dent is a partial ground of the sphere's being roughly spherical without making the difference between its being so and its not being so. Difference-making is the component that grounding shares with causation, and it is the component that fails to be transitive in both; the in-virtue-of relation that remains when it is set aside is transitive in grounding, and has no transitive counterpart in causation.

The analogy therefore supports the unity of grounding only in its weakest form, by showing that a relation can be one while operating through many mechanisms. The unity theorem does not depend on it. The theorem rests on two features of mixed chains, and causation has the first, since causal chains do not return to their origin, but lacks the second. Grounding is one relation; causation is another; and they are alike in being one relation each.

\begin{thebibliography}{99}\small
\bibitem[Bennett(2017)]{Bennett2017} Bennett, K. (2017). \emph{Making Things Up}. Oxford University Press.
\bibitem[Hall(2000)]{Hall2000} Hall, N. (2000). Causation and the price of transitivity. \emph{The Journal of Philosophy}, 97(4), 198--222.
\bibitem[Lewis(1973)]{Lewis1973} Lewis, D. (1973). Causation. \emph{The Journal of Philosophy}, 70(17), 556--567.
\bibitem[McDermott(1995)]{McDermott1995} McDermott, M. (1995). Redundant causation. \emph{The British Journal for the Philosophy of Science}, 46(4), 523--544.
\bibitem[Schaffer(2012)]{Schaffer2012} Schaffer, J. (2012). Grounding, transitivity, and contrastivity. In F.~Correia and B.~Schnieder (Eds.), \emph{Metaphysical Grounding: Understanding the Structure of Reality} (pp.~122--138). Cambridge University Press.
\bibitem[Schaffer(2016)]{Schaffer2016} Schaffer, J. (2016). Grounding in the image of causation. \emph{Philosophical Studies}, 173(1), 49--100.
\bibitem[Wilson(2018)]{Wilson2018} Wilson, A. (2018). Metaphysical causation. \emph{No\^us}, 52(4), 723--751.
\end{thebibliography}

\end{document}
